\section{Training Strategy}

\subsection{Optimization}

All stages use AdamW with $\beta = (0.9, 0.95)$ and weight decay 0.1.
Learning rates use cosine decay with linear warmup:

\begin{table}[H]
\centering\small
\caption{Training hyperparameters per stage.}
\begin{tabular}{ccccccc}
\toprule
\textbf{Stage} & \textbf{LR} & \textbf{Warmup} & \textbf{Batch} &
\textbf{Seq} & \textbf{Max Tokens} & \textbf{Patience} \\
\midrule
2 & $1.5 \times 10^{-4}$ & 1000 & 24 & 384 & 300M & 15 \\
3 & $1.5 \times 10^{-4}$ & 1000 & 24 & 384 & 300M & 15 \\
4 & $1.0 \times 10^{-4}$ & 1000 & 16 & 512 & 400M & 15 \\
5 & $1.0 \times 10^{-4}$ & 1000 & 16 & 512 & 500M & 15 \\
6 & $8.0 \times 10^{-5}$ & 1500 &  8 & 768 & 600M & 15 \\
\bottomrule
\end{tabular}
\end{table}

\textbf{Differential learning rates.}
At expansion stages (2, 4, 6), pretrained parameters receive a lower learning
rate (\texttt{pretrained\_lr\_mult} = 0.5) while newly initialized parameters
(cloned layers, widened FFN columns) use the full rate.  This prevents
catastrophic unlearning of pretrained weights during early training.

\subsection{Early Stopping}

Training terminates via a \emph{relative plateau detector}.  The target
validation metric (determined by \texttt{val\_key}) is monitored with patience
of 15 evaluation intervals.  A relative minimum delta of 0.1\% is used, along
with a minimum token count before exit is allowed (20--30M tokens depending on
stage).  All stages in this run exited via \texttt{plateau}.

\subsection{Gradient Clipping}

Global gradient norm is clipped to 1.0.  Both global and deep-layer gradient
norms are logged to detect instability during expansion.
